\section{Fixed-count transport}\label{sec:transport}

\subsection{The circular stars}

\begin{definition}[stars and the bow-tie]\label{def:star}
For $r\geq1$ put $N=2r+1$, $u_k=(\cos\frac{2\pi k}N,\sin\frac{2\pi k}N)$ for
$k\in\ZZ/N$, and
\[
K_r=(\mu_1,\ldots,\mu_N),\qquad \mu_t=u_{r(t-1)},
\]
the polygon visiting every $r$-th point of the regular $N$-gon in turn. Put
$K_{-r}=\overline{K_r}$. The \emph{bow-tie} is
$K_0=\bigl((0,0),(2,2),(0,2),(2,0)\bigr)$.
\end{definition}

\begin{lemma}[the stars are generic]\label{lem:star-generic}
Let $r\geq1$ and $N=2r+1$. Then $\gcd(r,N)=1$, and the following hold.
\begin{enumerate}
\item[(i)] $\sigma K_r=R(K_r)$, where $R$ is the rotation of the plane by
$2\pi r/N$ about the origin. Consequently every principal turn of $K_r$
equals $2\pi r/N\in(0,\pi)$, all turns are left, and $\rot(K_r)=r$.
\item[(ii)] Every edge line of $K_r$ is tangent to the circle of radius
$\rho=\cos(\pi r/N)$ about the origin, at the midpoint of the edge, and the
origin lies strictly to the left of every directed edge.
\item[(iii)] $K_r$ is generic; hence $K_{-r}$ is generic with all turns
right and $\rot(K_{-r})=-r$.
\item[(iv)] $K_0$ is generic, with
$(\chi_{123},\chi_{124},\chi_{134},\chi_{234})=(+1,-1,-1,+1)$, turns
$(\tau_1,\tau_2,\tau_3,\tau_4)=(-1,+1,+1,-1)$, $\rot(K_0)=0$ and
$X(K_0)=\{\{1,3\}\}$.
\end{enumerate}
\status{new (round~4: $\gcd(r,2r+1)=1$ stated, proved and cited at its two uses; adopted from the Codex proposal sm4\_round4\_triple\_star\_v1 (seal 9a82b59a\ldots), F125, re-read by the author; F-25-125; cleared by bench~A on SM5 (2026-09-06T00:32:58Z))}
\end{lemma}
\begin{proof}
First, every common positive divisor of $r$ and $N$ divides
$N-2r=1$, so $\gcd(r,N)=1$.
For integer indices $j,k$, if $rj\equiv rk\pmod N$ then
$N$ divides $r(j-k)$. Multiplying that divisibility by $2$ and
using $2r=N-1$ shows that $N$ divides $j-k$.
Thus multiplication by $r$ permutes the residue classes modulo $N$,
so the displayed star vertices are distinct.
Moreover $R^k$ is the identity precisely when $kr/N$ is an integer.
The same divisibility argument makes this equivalent to $N\mid k$;
the least positive such $k$ is $N$, proving that $R$ has order $N$.

(i) $(\sigma K_r)_t=\mu_{t+1}=u_{rt}=R(u_{r(t-1)})=R(\mu_t)$. Hence
$\lt_t=R(\lt_{t-1})$ for every $t$, so each principal turn is the angle of
$R$ reduced to $(-\pi,\pi)$, which is $2\pi r/N$ itself because $2r<N$; the
$N$ turns sum to $2\pi r$. (ii) $E_t$ is the chord from $u_a$ to $u_{a+r}$
with $a=r(t-1)$. The counterclockwise arc from $u_a$ to $u_{a+r}$ has angle
$2\pi r/N<\pi$, so the midpoint of the chord is at distance $\cos(\pi r/N)$
from the origin and the chord line is tangent there to the circle of that
radius; a chord directed along a counterclockwise arc of angle less than
$\pi$ has the centre on its left. (iii) (G1): the residue permutation proved above gives $N$ distinct
vertices on the unit circle, and no three points of a circle are collinear.
(G2): the $N$ edge lines are pairwise distinct tangent lines of the circle of
radius $\rho$ (their tangency points are the $N$ distinct images of one
midpoint under the powers of $R$, whose order was proved to be $N$). Through a point
outside a circle pass exactly two tangent lines of the circle, so no three of
the edge lines are concurrent, and a fortiori no three edge interiors are.
The statements for $K_{-r}$ are Lemma~\ref{lem:shift}(ii)--(iv). (iv) The
four determinants are $\det((2,2),(0,2))=4$, $\det((2,2),(2,0))=-4$,
$\det((0,2),(2,0))=-4$, $\det((-2,0),(0,-2))=4$; the turns are
$\tau_1=\chi_{412}=\chi_{124}$, $\tau_2=\chi_{123}$, $\tau_3=\chi_{234}$,
$\tau_4=\chi_{341}=\chi_{134}$ (Lemma~\ref{lem:chi-basic}(i)). The edges
$E_1=[(0,0),(2,2)]$ and $E_3=[(0,2),(2,0)]$ meet at $(1,1)$, while $E_2$ and
$E_4$ are disjoint parallel segments; there is one crossing point, so (G2)
holds. The principal turns are $-\frac{3\pi}4,\frac{3\pi}4,\frac{3\pi}4,-\frac{3\pi}4$,
which sum to $0$.
\end{proof}

\subsection{Connectivity of the regular fibres}

% Printed in round 1 from the Codex pack candidate_transport_v1 (C021); no historical machine check is a premise.
\begin{lemma}[positive closing lengths along a direction path]
\label{lem:transport-lengths}
Let $u_1(t),\ldots,u_n(t)$ be continuous unit vectors on a compact interval,
lying in no closed semicircle at any parameter. There are continuous strictly
positive lengths $l_i(t)$ with $\sum_i l_i(t)u_i(t)=0$. At either endpoint,
any prescribed positive closing lengths can be joined to the selected lengths
while keeping the directions fixed.
\status{proved (refereed: bench A, 2026-09-05T15:51:04Z; transcribed from CV thm:mycyclic, prose proof in d8\_transport (C021))}
\end{lemma}
\begin{proof}
For finitely many planar unit vectors, lying in no closed semicircle is
equivalent to the origin being in the interior of their convex hull.
Indeed a supporting line through the origin for a boundary point of that
convex polygon, or a separating line if the origin is outside it, puts all
vectors in one closed half-plane through the origin. Conversely, a closed
half-plane through the origin cannot contain a neighbourhood of the origin.
This is also valid for a segment or a point as convex hull, since either is
contained in a closed half-plane through the origin whenever the origin is
not an interior point in the plane.

Fix a parameter and put $b=\sum_i u_i$. Interiority gives an $\epsilon>0$
such that $-\epsilon b$ belongs to the convex hull. Write
$-\epsilon b=\sum_i a_i u_i$ with $a_i\geq0$. Then
\begin{equation}\label{transport:positive-close}
 \sum_i(a_i+\epsilon)u_i=-\epsilon b+\epsilon\sum_i u_i=0,
 \qquad a_i+\epsilon>0.
\end{equation}
Thus there is a strictly positive closing vector at each parameter.

Fix such a vector $l$ at $t_0$. Choose $p,q$ with
$\det(u_p(t_0),u_q(t_0))\ne0$, possible because the vectors are not all
parallel. Let $E(t)=\sum_i l_i u_i(t)$ and
$D(t)=\det(u_p(t),u_q(t))$. Correct just two coordinates by
\begin{equation}\label{transport:cramer}
 \delta_p(t)=-\frac{\det(E(t),u_q(t))}{D(t)},\qquad
 \delta_q(t)=-\frac{\det(u_p(t),E(t))}{D(t)}.
\end{equation}
The determinant expansion of a vector in the basis $(u_p,u_q)$ gives
$E+\delta_pu_p+\delta_qu_q=0$. The correction is continuous and vanishes
at $t_0$, so all corrected lengths remain positive on a neighbourhood of
$t_0$. This constructs local continuous positive closing vectors.

For a one-point parameter interval the pointwise choice already suffices.
Otherwise choose a finite cover by proper smaller relative intervals, with
nonempty complements, whose closures lie in these neighbourhoods.
Continuous nonnegative weights subordinate to the smaller
intervals, summing to one, can be obtained by taking distances to their
complements and normalizing their sum. Multiply each weight by the local
closing vector on its neighbourhood and sum. The products extend continuously
by zero outside the smaller intervals because the local vectors are bounded
on their closures. The result is a continuous positive closing vector: closure
is linear and every positively weighted local vector has all entries positive.
Finally, for fixed directions the set of strictly positive closing lengths is
convex. Straight segments in that set join any prescribed endpoint lengths
to the selected ones.
\end{proof}

\begin{lemma}[lifting a semicircle to one real interval]
\label{lem:transport-angle-interval}
If a finite real sequence $\theta_0,\ldots,\theta_N$ has
$|\theta_{i+1}-\theta_i|<\pi$, and all its unit directions lie in a closed
semicircle, then all its real entries lie in one interval of length $\pi$.
Conversely, containment of all entries in such an interval puts their unit
directions in a closed semicircle.
\status{proved (refereed: bench A, 2026-09-05T15:51:04Z; transcribed from CV thm:mycyclic, prose proof in d8\_transport (C021))}
\end{lemma}
\begin{proof}
The inverse image of a fixed closed semicircle under
$\theta\mapsto(\cos\theta,\sin\theta)$ is a union of intervals
$[\alpha+2\pi k,\alpha+\pi+2\pi k]$. Distinct consecutive intervals have
a gap of length $\pi$. A step of absolute size strictly less than $\pi$
cannot move from one interval to another, so induction fixes one interval
for the whole sequence. The converse follows by projection of that interval.
\end{proof}

\begin{theorem}[connectivity of the fibres]\label{thm:mycyclic}
Let $(n,r)$ be admissible and let $P,P'\in\mathcal R_n$ be labelled tuples
with $\rot(P)=\rot(P')=r$. If $(n,r)\ne(4,0)$, a continuous path in
$\mathcal R_n\cap\rot^{-1}(r)$ joins $P$ to $P'$. If $(n,r)=(4,0)$,
there is a $k\in\mathbb Z/4$ and such a path from $P$ to $\sigma^kP'$.
The four shifts of $K_0$ lie in the four distinct labelled components,
distinguished by their turn words. Thus every fibre of cyclic polygon orbits
is path connected.
\status{proved (refereed: bench A, 2026-09-05T19:29:31Z; transcribed from CV thm:mycyclic --- the $k=0$ refinement from CV's proof, not its statement --- with RC root:four-components for the realizability of the four words at $(4,0)$; the two locators compose, each supplying a half (C021; F-25-49, F-25-104, F-25-129))}
\end{theorem}
\begin{proof}
Use tail directions $d_i=\mu_{i+1}-\mu_i$, unit directions $u_i=d_i/|d_i|$,
and principal turns $\vartheta_i\in(-\pi,\pi)$ from $u_{i-1}$ to $u_i$.
Multiplication of $u_i=e^{\mathrm i\vartheta_i}u_{i-1}$ around the cycle
shows that $\sum_i\vartheta_i$ is an integer multiple of $2\pi$.
Principal turns vary continuously on the regular locus, so their normalized
sum is constant on every path there. The symbols $\tau_i$ remain the signs
of these real principal turns, not the angles themselves.

\emph{Reconstructing a polygon from directions.}
A direction path with no antiparallel successive pair and with no closed
semicircle containment can be closed using
Lemma~\ref{lem:transport-lengths}. Starting at any continuous base point,
successively add the positive vectors $l_i(t)u_i(t)$; the final sum is zero,
so this gives a continuous closed labelled $n$-tuple. Every edge is nonzero,
and every successive pair has the prescribed principal turn in $(-\pi,\pi)$.
Consequently the tuple stays in $\mathcal R_n$. Interpolate the base point
between those of $P,P'$ and use the endpoint length segments of that lemma
to recover the exact endpoint tuples. No vertex is inserted or deleted.

\emph{Nonzero rotation.}
Choose real arguments $\beta,\beta'$ for the first unit directions and
interpolate these arguments and all principal turns linearly. The interpolated
turns lie strictly between $-\pi$ and $\pi$ and have sum $2\pi r$.
The lifted edge arguments, starting with $\theta_0(t)=\beta(t)$, are
\begin{equation}\label{transport:prefix}
 \theta_k(t)=\beta(t)+\sum_{m=1}^{k}\vartheta_{m+1}(t),
 \quad 0\leq k\leq n,
 \qquad \theta_n(t)-\theta_0(t)=2\pi r,
\end{equation}
where turn indices are cyclic. The first $n$ entries give the unit directions;
the last repeats the first direction because $r$ is an integer. If these
directions lay in a closed semicircle,
Lemma~\ref{lem:transport-angle-interval} would put the whole lifted sequence
in one interval of length $\pi$, implying $|2\pi r|\leq\pi$.
This contradicts $r\ne0$. The reconstruction above therefore joins the exact
labelled endpoints, with no cyclic shift.

\emph{The zero-rotation angle domain.}
For $r=0$, the lifted angles close as real numbers and form a cyclic vector
$\theta=(\theta_1,\ldots,\theta_n)$ satisfying
\begin{equation}\label{transport:zero-angle-domain}
 |\theta_{i+1}-\theta_i|<\pi\quad\text{cyclically}.
\end{equation}
The directions of any regular closed polygon lie in no closed semicircle.
For if a normal $v\ne0$ had $\langle v,d_i\rangle\geq0$ for every edge,
closure would force every scalar product to be zero. All directions would
then be on the two rays of one line. Antiparallel adjacent directions are
forbidden, so every edge would point along the same ray. Their positive
length sum could not be zero. By
Lemma~\ref{lem:transport-angle-interval}, the cyclic angle vector of a
zero-rotation polygon consequently has
$\max_i\theta_i-\min_i\theta_i>\pi$.

For an ordered pair of indices with forward cyclic distance
$d\in\{2,\ldots,n-2\}$, define
\begin{equation}\label{transport:chart}
 \mathcal A(h,d)=\{\theta:\eqref{transport:zero-angle-domain}\text{ holds},
                  \ \theta_h-\theta_{h+d}>\pi\}.
\end{equation}
These open convex charts cover all zero-rotation endpoints: take a maximum
and a minimum; they cannot be cyclically adjacent by
\eqref{transport:zero-angle-domain}. Conversely every chart point has
directions lying in no closed semicircle, by
Lemma~\ref{lem:transport-angle-interval}. Every continuous angle path inside
the union can thus be reconstructed as above.

\emph{Connected chart overlaps for $n\geq5$.}
For $2\leq d\leq n-3$, both of the following chart pairs intersect:
\begin{equation}\label{transport:chart-moves}
 \mathcal A(h,d),\mathcal A(h,d+1);
 \qquad \mathcal A(h,d+1),\mathcal A(h+1,d).
\end{equation}
Here is an explicit witness, reading indices $s=0,\ldots,n-1$ from $h$.
Set the high value to $3\pi/2$ and the low value to zero. For the first pair,
put the high value at $s=0$, descend linearly to zero at $s=d$, keep zero
through $s=d+1$, and rise linearly back to the high value at $s=n$.
For the second pair, keep the high value at $s=0,1$, descend linearly to zero
at $s=d+1$, and rise linearly back to the high value at $s=n$.
Every ramp has at least two steps, so each step has magnitude at most
$3\pi/4<\pi$. Each specified high-minus-low difference is $3\pi/2>\pi$,
proving both overlaps. The cyclic last step is one of these ramp steps.

The first moves reduce any distance to two. At distance two, the two moves
$\mathcal A(h,2)\to\mathcal A(h,3)\to\mathcal A(h+1,2)$ advance the high
index by one, and are valid even for $n=5$. Hence their overlap graph is
connected. Join endpoint angle vectors to consecutive overlap witnesses by
straight segments in their convex charts. The resulting finite path of
feasible angle vectors reconstructs a regular path from $P$ to $P'$, again
with no relabelling.

\emph{The exceptional four-vertex fibre.}
For $n=4$ the charts are just $\mathcal A(h,2)$ for four indices $h$.
They are nonempty: values $3\pi/2,3\pi/4,0,3\pi/4$, beginning at $h$,
satisfy all their inequalities. If $\theta_h-\theta_{h+2}>\pi$, the two
intermediate entries satisfy
\begin{equation}\label{transport:four-orders}
 \theta_h>\theta_{h+1}>\theta_{h+2},\qquad
 \theta_h>\theta_{h+3}>\theta_{h+2}.
\end{equation}
For example $\theta_{h+1}>\theta_h-\pi>\theta_{h+2}$, and
$\theta_{h+1}<\theta_{h+2}+\pi<\theta_h$; the other case is identical.
Thus the turn signs, beginning at vertex $h$, are $(+,-,-,+)$.
The four charts force the four distinct cyclic rotations of this word.
They are therefore disjoint, and every rotation-zero four-tuple belongs to
one of them. In particular no such tuple has a zero turn. Along a regular
zero-rotation path all four nonzero turn signs remain fixed by continuity,
so different words cannot be joined. Two tuples with the same word belong
to the same convex angle chart; their real argument choices can differ by
a common multiple of $2\pi$, which changes neither the chart inequalities
nor the reconstruction. Interpolation inside that chart and the positive-length
construction join them. There are exactly four components.

For $K_0=((0,0),(2,2),(0,2),(2,0))$ the corner determinants are
$(-4,4,4,-4)$. Its principal turns are the corresponding signs times
$3\pi/4$, so its rotation is zero. Cyclic shifts rotate this sign word and
exhibit all four components. Shifting $P'$ to match the word of $P$ gives
the asserted path in the exceptional case.

Finally a regular triangle is noncollinear: three collinear nonzero closing
edges must reverse direction at a vertex. All three turn determinants of a
noncollinear triangle equal its oriented area determinant, so the three
principal turns have one sign. Their sum is a nonzero multiple of $2\pi$
of absolute value less than $3\pi$; its rotation is $+1$ or $-1$.
There is no zero-rotation case at $n=3$. These cases exhaust admissible
arities. Passing any constructed path to the cyclic quotient identifies
its shifted target with the original target orbit, proving the orbit statement.
\end{proof}

\paragraph{Collision boundary of this statement.}
The regular locus forbids zero edges and antiparallel consecutive directions;
it does not forbid equality of nonadjacent vertices. Such equality may already
occur at the endpoints above, so no collision-free strengthening is asserted
for arbitrary regular endpoints.

\subsection{Transport through simple walls}



\begin{lemma}[transport]\label{lem:transport}
Let $P,Z$ be generic labelled tuples in the same fibre $(n,r)$. There are
$k\in\ZZ/n$, with $k=0$ unless $(n,r)=(4,0)$, and a piecewise-affine path in
$\mathcal R_n$ from $P$ to $\sigma^kZ$ that is generic except at finitely
many parameters, at each of which it is a simple wall germ of type (F), (V),
(T), (E) or (C); at each (F) the deletion is generic, and at each (V) both
halves are generic, all with fewer than $n$ vertices. \status{proved (refereed: bench A, 2026-09-05T18:22:57Z; transcribed from SM1 lem:transport, from Theorems~\ref{thm:mycyclic}, \ref{thm:relgp} and Lemma~\ref{lem:children}; proof paraphrased in the base (C021); F-25-76)}
\end{lemma}
\begin{proof}
Theorem~\ref{thm:mycyclic} supplies the path from $P$ to $\sigma^kZ$
in $\mathcal R_n$. Its rotation is constant by Lemma~\ref{lem:rot}(ii).
Theorem~\ref{thm:relgp} replaces it by a path with only the five listed
simple wall types. Lemma~\ref{lem:children} gives generic deletions and
halves.
\end{proof}

\subsection{Anchors}

\begin{lemma}[rotation of a soft insertion]\label{lem:soft-rotation}
Let $P$ be generic, $q$ admissible at $j$, and $\varepsilon$ small
(Lemma~\ref{lem:soft-generic}). Then $\rot(P_\varepsilon)=\rot(P)$ in the
same-sign and mixed sectors, and $\rot(P_\varepsilon)=\rot(P)-\tau_j(P)$ in
the loop sector. Moreover the same-sign sector is the open cone of directions
strictly between $\lt_{j-1}$ and $\lt_j$ (the wedge of the turn), the loop
sector is the open cone strictly between $-\lt_{j-1}$ and $-\lt_j$, and the
mixed sector is the complement of the closures of these two cones.
\status{new}
\end{lemma}
\begin{proof}
Put $u=\lt_{j-1}$, $v=\lt_j$, and
$\theta=\operatorname{Arg}(v/u)\in(-\pi,\pi)$, using the complex
identification of the oriented plane. Let
$\alpha=\operatorname{Arg}(q/u)$ and
$\beta_0=\operatorname{Arg}(v/q)$ be principal turns. All three
exist because admissibility makes the relevant pairs nonparallel.
At the new point the actual turn is
$\beta_\varepsilon=\operatorname{Arg}((v-\varepsilon q)/q)$,
which converges to $\beta_0$. The turn at the successor vertex also
changes, but converges to its old value because its incoming direction
$v-\varepsilon q$ tends to $v$ and the old corner is regular.
All other original turns except the replaced turn at $j$ are unchanged.
By Lemmas~\ref{lem:soft-generic} and~\ref{lem:rot}(ii), the rotation
of $P_\varepsilon$ is one fixed integer on a sufficiently small
positive interval. Taking the limit of the full sum of turns therefore
gives the exact integer identity
\begin{equation}\label{eq:soft-rotation-limit}
2\pi\bigl(\rot(P_\varepsilon)-\rot(P)\bigr)
=\alpha+\beta_0-\theta.
\end{equation}
The left side is constant on that interval, not an assertion that the
successor turn is unchanged at finite $\varepsilon$.

The three direction ratios multiply as $(q/u)(v/q)=v/u$, so
$\alpha+\beta_0-\theta$ is a multiple of $2\pi$.
Moreover $\sgn\alpha=-\chi_-$ and $\sgn\beta_0=-\chi_+$.
In the same-sign sector both have the sign $\tau_j=\sgn\theta$.
Their sum and $\theta$ lie in the same open interval of length $2\pi$,
so the multiple is zero. In the mixed sector $\alpha,\beta_0$ have
opposite signs, making their sum lie in $(-\pi,\pi)$; the difference
from $\theta$ is again strictly between $-2\pi$ and $2\pi$, so is zero.
In the loop sector both signs are $-\tau_j$. When $\tau_j=1$,
the difference lies in $(-3\pi,0)$ and must be $-2\pi$; when
$\tau_j=-1$ it lies in $(0,3\pi)$ and must be $2\pi$.
Equation~\eqref{eq:soft-rotation-limit} gives exactly the stated rotations.

For the cone description, set the angle of $u$ equal to zero and that
of $v$ equal to $\theta$. The determinant signs change only at the
four rays $u,v,-u,-v$. The same-sign inequalities put $q$ strictly
inside the turn wedge from $u$ to $v$ of angular width $|\theta|<\pi$;
the loop inequalities put it inside the opposite wedge. The remaining
two open angular intervals are exactly the mixed sector. This proves
the cone statements as well.
\end{proof}

\begin{definition}[anchors]\label{def:anchors}
Let $(n,r)$ be admissible with $n\geq4$.
\begin{enumerate}
\item[(Z)] If $(n-1,r)$ is admissible, a \emph{zero anchor} is a generic
$n$-gon of the form $P^{(j,q)}_\varepsilon$ (Definition~\ref{def:soft}) in
the mixed sector, for some generic $(n-1)$-gon $P$ of rotation $r$, some
vertex $j$ of $P$, some admissible $q$ and some small $\varepsilon>0$.
\item[(L)] If $(n,r)$ is minimal with $|r|\geq2$, a \emph{loop anchor} is a
generic $n$-gon of the form $P^{(j,q)}_\varepsilon$ in the loop sector, for
some generic $(n-1)$-gon $P$ of rotation $r-\sgn(r)$ and some vertex $j$ of
$P$ with $\tau_j(P)=-\sgn(r)$.
\item[(L$_0$)] If $(n,r)=(4,0)$, a \emph{loop anchor} is a generic $4$-gon
of the form $P^{(j,q)}_\varepsilon$ in the loop sector for a triangle $P$
and any of its vertices $j$.
\end{enumerate}
In all three cases $q$ is admissible and $0<\varepsilon<\varepsilon_0$, where
$\varepsilon_0=\varepsilon_0(P,j,q)>0$ is the bound of
Lemma~\ref{lem:soft-generic}: on $(0,\varepsilon_0)$ the insertion is generic
and lies in one chamber, so that Lemma~\ref{lem:soft-rotation} gives its
rotation throughout. This geometric bound is part of the anchor data; it is
not a bound uniform over the functions to which a soft theorem is applied.
In each case $P$ is the \emph{parent} of the anchor, $j$ its \emph{insertion
vertex}, and the soft edge of the anchor is its edge $E_j$ (from $\mu_j$ to
$\mu_*=\mu_{j+1}$ in the labelling of the anchor).
\end{definition}

\begin{proposition}[anchors exist]\label{prop:anchors-exist}
For every admissible $(n,r)$ with $n\geq4$, exactly one of the cases
(Z), (L), (L$_0$) of Definition~\ref{def:anchors} applies. An anchor
of that type exists, is generic, has $n$ vertices and rotation $r$. For a given generic parent $P$ and vertex $j$, the mixed
sector and the loop sector both contain admissible vectors $q$, so in case
(Z) the insertion vertex may be prescribed. \status{new (round~4: the exhaustion of the three cases printed, adopted from the Codex proposal sm4\_round4\_core\_v1, P-25-23, F118, re-read by the author; F-25-118; cleared by bench~A on SM5 (2026-09-06T00:32:58Z) before the round-5 disjointness sentence of the exhaustion paragraph (F-25-150) and the citations of Proposition~\ref{prop:chambers} and of the density clause of Lemma~\ref{lem:fibres} in the loop-anchor construction (F-25-146))}
\end{proposition}
\begin{proof}
\emph{Exhaustion.} If $(n-1,r)$ is admissible, case (Z) applies.
Otherwise $n-1\geq3$, so Definition~\ref{def:admissible} says that
either $2|r|\geq n-1$ or $(n-1,r)=(3,0)$. In the first case,
$2|r|<n$ and integrality give $2|r|=n-1$. Thus
$n=2|r|+1$, and $n\geq4$ forces $|r|\geq2$, giving (L).
The exceptional predecessor gives exactly $(n,r)=(4,0)$, case
(L$_0$). These alternatives are disjoint: in case (L) the predecessor
$(n-1,r)=(2|r|,r)$ has $2|r|=n-1$ and in case (L$_0$) it is $(3,0)$, so in
neither is it admissible (Definition~\ref{def:admissible}), which excludes
(Z); and (L) has $|r|\geq2$ while (L$_0$) has $r=0$. In particular $(4,1)$
and $(4,-1)$ have admissible triangular predecessors and belong
 to (Z), not (L).

\emph{Sectors.} By Lemma~\ref{lem:soft-rotation} each sector is a nonempty
open cone of directions; admissibility removes finitely many lines
($\det(q,\mu_k-\mu_j)=0$, $k\neq j$), so each sector contains admissible
vectors. For such $q$ and $0<\varepsilon<\varepsilon_0$, as Definition~\ref{def:anchors}
requires, Lemma~\ref{lem:soft-generic} makes $P_\varepsilon$ generic and
Lemma~\ref{lem:soft-rotation} gives its rotation.

(Z) A generic $(n-1)$-gon of rotation $r$ exists by Lemma~\ref{lem:fibres};
in the mixed sector the rotation is preserved.

(L) Let $r\geq2$. The star $K_{r-1}$ is generic with all turns left and
rotation $r-1$ (Lemma~\ref{lem:star-generic}; for $r=2$ it is the
counterclockwise triangle). Subdivide one edge at its midpoint, obtaining a
$2r$-gon in $\mathcal R_{2r}$ with one zero turn and rotation $r-1$
(Lemma~\ref{lem:rot}(iii)). Write the subdivided edge as $B-A=w$
and its moved midpoint as $M=(A+B)/2+\delta v$, where $\delta>0$
and $\det(w,v)>0$: the displacement is to the left of the directed edge.
The new corner determinant is
\begin{equation}\label{eq:anchor-midpoint-turn}
\det(M-A,B-M)=-\delta\det(w,v)<0.
\end{equation}
Thus its turn is right. The two neighbouring turn determinants are
strictly positive at $\delta=0$ and remain positive for a short move.
All edge lengths remain positive; the initially positive-flat new
corner and all other regular corners stay in the open regular locus.
Rotation therefore remains $r-1$ by Lemma~\ref{lem:rot}(ii). The
generic locus is open (Proposition~\ref{prop:chambers}) and dense (the
density clause of Lemma~\ref{lem:fibres}), so a
small perturbation gives a generic $2r$-gon $P$ of rotation $r-1$ with a
vertex $j$ of turn $\tau_j=-1=-\sgn(r)$. A loop-sector insertion at $j$ gives
a generic $(2r+1)$-gon of rotation $r-1-\tau_j=r$. For $r\leq-2$ apply the
reflection $(x,y)\mapsto(x,-y)$ to the construction for $-r$: it negates every
chirotope and every rotation number, maps generic polygons to generic
polygons, and maps a loop-sector insertion to a loop-sector insertion (the
sectors are defined by sign patterns relative to $\tau_j$, which all flip
together).

(L$_0$) Insert a loop at any vertex of the counterclockwise triangle: the
rotation becomes $1-1=0$, and the result is generic with four vertices.
\end{proof}

\begin{proposition}[anchor values]\label{prop:anchor-values}
Let $F$ be a function on generic polygons of all arities, constant
on chambers and satisfying the soft theorem in the form of
Theorem~\ref{thm:C-soft} at every admissible soft insertion into
a generic polygon. With $P$ the parent and $j$ the insertion vertex,
\[
F(Z)=0\quad\text{for every zero anchor},\qquad
F(Y)=\tau_j(P)F(P)\quad\text{for every loop anchor}.
\]
In case (L), $\tau_j(P)=-\sgn(r)$; in case (L$_0$), it is the
orientation sign of the parent triangle. The function $C$ satisfies
these hypotheses. The same identities hold for $A_g$ at every
anchor root other than the soft edge, with the corresponding
parent root used on the right-hand side.
\status{new (round~4: chamber constancy added to the hypothesis and the function-dependent soft threshold handled by transport along the initial chamber; adopted from the Codex proposal sm4\_round4\_core\_v1, P-25-23, F119, re-read by the author; F-25-119; cleared by bench~A on SM5 (2026-09-06T00:32:58Z))}
\end{proposition}
\begin{proof}
Fix an anchor with parameter $\varepsilon$ and geometric bound
$\varepsilon_0$ from Definition~\ref{def:anchors}. For its fixed
parent, vertex and admissible vector, the soft hypothesis for $F$
supplies a bound $\delta_F>0$. Choose
$0<\varepsilon'<\min\{\varepsilon,\delta_F\}$.
The whole parameter interval between $\varepsilon'$ and
$\varepsilon$ is contained in $(0,\varepsilon_0)$, so both tuples
lie in the same chamber by Lemma~\ref{lem:soft-generic}.
Chamber constancy equates their $F$ values. At $\varepsilon'$ the
soft theorem applies. Its multiplier is zero in the mixed sector,
since the attachment signs are opposite. In the loop sector both
attachment signs equal $\tau_j(P)$, so the multiplier is
$\tau_j(P)$. This proves both identities at the original
$\varepsilon$, without requiring a bound uniform over functions.

Proposition~\ref{prop:C-chamber} and Theorem~\ref{thm:C-soft}
supply the two hypotheses for $C$. For $A_g$, fix the physical
nonsoft edge throughout the labelled family and its corresponding
edge of the parent. Theorem~\ref{thm:A-soft} gives the required
identity on an initial interval for that root. Throughout the
geometric interval the tuple remains generic, so every initially
nonzero determinant keeps its sign along the connected parameter
interval. Proposition~\ref{prop:A-chamber} therefore makes $A_g$
constant along that family. The same smaller-parameter argument
proves the rooted assertion, without using root independence.
\end{proof}
